\subsubsection{Marco Conceptual}

\begin{itemize}
    \item \textbf{SaaS (Software as a Service).} Modelo de distribución de software en el que las aplicaciones se alojan en servidores de nube y se ofrecen a través de internet bajo suscripción, sin instalación local (\cite{mellgrance2011}).

    \item \textbf{Bot conversacional / Chatbot.} Programa informático diseñado para simular conversaciones con personas mediante lenguaje natural o flujos de diálogo estructurados, procesando mensajes y generando respuestas automáticas (\cite{adamopoulou2020}).

    \item \textbf{Agente de autorespuesta.} Tipo de chatbot orientado a atención institucional o comercial, que responde consultas frecuentes de forma autónoma y escala a operador humano cuando es necesario.

    \item \textbf{API (Application Programming Interface).} Conjunto de definiciones y protocolos que habilita la comunicación e integración entre sistemas de software para intercambiar datos y funcionalidades (\cite{fielding2000}).

    \item \textbf{Baileys (WhatsApp Web API).} Biblioteca de JavaScript/TypeScript que permite integrar WhatsApp Web de forma programática para enviar, recibir y gestionar mensajes automatizados.

    \item \textbf{Telegram Bot API.} API pública y gratuita de Telegram para crear bots que interactúan con usuarios, procesan mensajes, envían multimedia y ejecutan comandos en la plataforma (\cite{telegram2024}).

    \item \textbf{PLN (Procesamiento de Lenguaje Natural).} Rama de la inteligencia artificial que estudia la interacción entre computadoras y lenguaje humano para comprender, interpretar y generar texto significativo (\cite{manningschutze1999}).

    \item \textbf{Omnicanalidad.} Estrategia que garantiza una experiencia de usuario coherente y continua, independientemente del canal utilizado, unificando historiales y criterios de atención (\cite{rigby2011}).

    \item \textbf{Multitenencia (Multitenancy).} Arquitectura donde una misma instancia de aplicación atiende múltiples clientes de forma aislada y segura, compartiendo eficientemente recursos de infraestructura (\cite{armbrust2010}).

    \item \textbf{Trazabilidad conversacional.} Capacidad del sistema para registrar y recuperar el historial completo de interacciones, facilitando auditorías, análisis de comportamiento y continuidad de atención.

    \item \textbf{Escalabilidad.} Propiedad del sistema para incrementar capacidad de procesamiento o usuarios atendidos conforme aumenta la carga, sin degradar rendimiento ni rediseñar la arquitectura.

    \item \textbf{Handoff (Transferencia de canal).} Proceso mediante el cual una conversación automatizada se transfiere a un operador humano, preservando contexto e historial para asegurar continuidad (\cite{folstadskjuve2019}).

    \item \textbf{CSAT (Customer Satisfaction Score).} Indicador de satisfacción obtenido mediante encuestas breves posteriores a la interacción, que mide valoración del usuario en escala numérica definida.

    \item \textbf{Tríada CIA (Confidencialidad, Integridad, Disponibilidad).} Marco de seguridad de la información que establece acceso autorizado, exactitud de datos y disponibilidad oportuna de la información (\cite{iso27001_2022}).
\end{itemize}